\documentclass[11pt,a4paper]{article}

% =========================================================
%  Style (arXiv template provided with the project)
% =========================================================
\usepackage{arxiv}

% =========================================================
%  Language, fonts, input encoding
% =========================================================
\usepackage[utf8]{inputenc}
\usepackage[T2A,T1]{fontenc}
\usepackage[english,russian,main=english]{babel}

% =========================================================
%  Core typography / math / graphics
% =========================================================
\usepackage{amsmath,amsfonts,amssymb}
\usepackage{mathtools}
\usepackage{nicefrac}
\usepackage{microtype}
\usepackage{graphicx}
\usepackage{booktabs}
\usepackage{multirow}
\usepackage{array}
\usepackage{makecell}
\usepackage{xcolor}
\usepackage{setspace}
\usepackage{enumitem}
\usepackage{caption}
\usepackage{subcaption}
\usepackage{natbib}
\usepackage{hyperref}
\usepackage{doi}
\usepackage{cleveref}
\usepackage{url}

% =========================================================
%  Geometry override to A4 (arxiv.sty loaded geometry as letterpaper)
% =========================================================
\AtBeginDocument{\newgeometry{a4paper,margin=2.4cm,headheight=14pt,headsep=20pt,footskip=24pt}}

% =========================================================
%  Hyperref / PDF metadata
% =========================================================
\hypersetup{
  colorlinks=true,
  linkcolor=black,
  citecolor=black,
  urlcolor=blue,
  pdftitle={Weakly Supervised Semantic Segmentation for Plant Disease Localization},
  pdfauthor={Denis Bakin},
  pdfkeywords={weakly supervised semantic segmentation, plant disease, class activation maps, Siamese network, localization probe}
}

% =========================================================
%  Title
% =========================================================
\title{Weakly Supervised Semantic Segmentation for Plant Disease Localization}
\author{Denis Bakin}
\date{April 2026}
\renewcommand{\shorttitle}{WSSS for Plant Disease Localization}

% =========================================================
%  Shortcuts
% =========================================================
\newcommand{\Leq}{\mathcal{L}_{\mathrm{eq}}}
\newcommand{\Lcon}{\mathcal{L}_{\mathrm{con}}}
\newcommand{\Ldist}{\mathcal{L}_{\mathrm{distill}}}
\newcommand{\Lcls}{\mathcal{L}_{\mathrm{cls}}}
\newcommand{\Ltotal}{\mathcal{L}_{\mathrm{total}}}
\newcommand{\IoU}{\mathrm{IoU}}
\newcommand{\mIoU}{\mathrm{mIoU}}

\begin{document}

% =========================================================
%  Title page (cover)
% =========================================================
% \begin{titlepage}
\newpage

{\setstretch{1.0}
\begin{center}
NATIONAL RESEARCH UNIVERSITY\\
HIGHER SCHOOL OF ECONOMICS
\\
\bigskip
Faculty of Computer Science\\
Bachelor’s Programme “Applied Mathematics and Information Science”
\end{center}
}

\vspace{7em}

\begin{center}
{\bf BACHELOR'S THESIS}\\
{\bf Research Project on the Topic:}\\
{\bf Weakly Supervised Semantic Segmentation for Plant Disease Localization}\\
\end{center}

\vspace{2em}

{\bf Submitted by the Student: \vspace{2mm}}

{\setstretch{1.1}
\begin{tabular}{l@{\hskip 1.5cm}l}
group \#\foreignlanguage{russian}{БПМИ}221, 4th year of study & 
Denis Bakin
\end{tabular}}

\vspace{1em}
{\bf Approved by the Supervisor: \vspace{2mm}}

{\setstretch{1.1}
\begin{tabular}{l@{\hskip 1.5cm}l}
(Supervisor name)\\
(Supervisor position)\\
Faculty of Computer Science, HSE University
\end{tabular}}

\vspace{\fill}

\begin{center}
Moscow 2026
\end{center}

\end{titlepage}


% =========================================================
%  Abstract (English)
% =========================================================
\begin{abstract}
\noindent
We study weakly-supervised semantic segmentation (WSSS) for plant disease
localization on the in-the-wild \textsc{PlantSeg} dataset, where only
image-level labels are available at training time. We reproduce three
strong PASCAL VOC baselines (MCTformer, PSA/Random-Walk, WeakCLIP) and
show that they transfer poorly to plant lesions: the best full-pipeline
disease IoU is $33.3\%$ against a fully-supervised upper bound of
$70.1\%$. As the main methodological contribution we introduce a
\emph{shallow-probe localization-capacity} analysis that upper-bounds the
pixel information latent in any intermediate layer, and use it to (i)~uncover
$42$--$62\%$ disease IoU hidden inside classifier-only SPDNet features and
(ii)~document a clean null result for equivariance and patch-contrastive
auxiliary losses.
\end{abstract}

\vspace{0.5em}
\keywords{weakly-supervised semantic segmentation \and plant disease \and
class activation maps \and Siamese networks \and localization probes \and
auxiliary losses}

% =========================================================
%  Abstract (Russian)  —  Аннотация
% =========================================================
\begin{otherlanguage}{russian}
\begin{abstract}
\noindent
В работе изучается задача слабо-управляемой семантической сегментации
(WSSS) для локализации болезней растений на полевом наборе
\textsc{PlantSeg}, где на этапе обучения доступны только метки уровня
изображения. Воспроизведены три базовых метода (MCTformer,
PSA/Random\,Walk, WeakCLIP); показано, что их прямой перенос из домена
PASCAL VOC даёт лишь $33{,}3\%$ IoU на классе disease при
полностью-управляемом верхнем пределе $70{,}1\%$. Основным
методологическим вкладом является \emph{анализ локализационной ёмкости
с помощью мелкого декодера-зонда}, позволяющий оценить локализационную
информацию, содержащуюся в произвольном слое сети. С его помощью
(i)~обнаружено $42$--$62\%$ IoU, скрытых в признаках классификатора
SPDNet, и (ii)~документирован отрицательный результат для
вспомогательных функций потерь $\Leq$ и $\Lcon$.
\end{abstract}
\end{otherlanguage}

% =========================================================
%  Table of contents
% =========================================================
\newpage
\pagenumbering{roman}
{\small\setstretch{1.05}\tableofcontents}
\newpage
\pagenumbering{arabic}
\setcounter{page}{1}

% =========================================================
%  1 — Introduction
% =========================================================
\section{Introduction}
\label{sec:intro}

Plant diseases cause recurring yield losses on the order of
$20$--$40\%$ of global crop production and directly affect food security,
supply chains and the economy~\citep{plantseg,qin_severity,garden}. Early,
accurate identification and spatial delineation of disease regions on the
plant is therefore a scientifically and economically valuable task.
Modern computer-vision pipelines can in principle provide such delineation
through \emph{semantic segmentation}, but pixel-accurate disease masks are
prohibitively expensive to produce: they require domain-expert annotators,
careful boundary labelling, and considerable time per image. The most
recently released PlantSeg dataset~\citep{plantseg}, with roughly
$7{,}700$ expertly annotated field images, is currently the largest public
pixel-labelled resource and illustrates how severely supply-constrained
the plant-disease segmentation domain is.

\paragraph{Research question.}
Given this scarcity of dense annotations, a natural research direction is
\emph{weakly-supervised semantic segmentation} (WSSS): learn a segmentation
model from image-level labels only (``the image contains disease $c$''),
leveraging the much cheaper annotations available in plant-phenotyping
collections such as PlantVillage~\citep{plantvillage}. WSSS has been
extensively studied on general-purpose benchmarks — notably PASCAL
VOC\,2012 — where state-of-the-art methods achieve mean IoU values on the
order of $60$--$75\%$ and are within a few percentage points of the
fully-supervised upper bound \citep{sec,affinitynet,weaktr,weakclip}. These
methods typically follow a common pattern: (i)~train an image-level
classifier, (ii)~extract \emph{class activation maps}
(CAMs)~\citep{cam,gradcam}, (iii)~refine the CAMs with dense
CRF~\citep{densecrf} and pixel-affinity~\citep{affinitynet} modules to
produce pseudo-masks, and (iv)~train a segmentation network on these
pseudo-masks.

The main claim of this thesis is that \emph{this recipe transfers poorly
to the plant-disease domain as it currently stands}, and that the gap
between PASCAL VOC WSSS ($\mIoU\approx65\%$) and plant-disease WSSS
(our best end-to-end result: $46\%$ $\mIoU$, $33\%$ disease IoU) is not
an artefact of hyperparameter tuning but a structural property of the
problem: plant lesions are small, sparse, low-contrast, and visually
similar across the background; moreover, a standard classifier trained
for disease \emph{class} prediction is not required to produce
spatially discriminative attention to solve its training objective.

\paragraph{Contributions.}
The thesis makes four main contributions, all building on the PlantSeg
and PlantVillage datasets:

\begin{enumerate}[label=(\roman*),leftmargin=*,itemsep=2pt]
\item A reproducible baseline evaluation of three well-known WSSS methods
(MCTformer~\citep{mctformer}, PSA\,+\,Random\,Walk~\citep{affinitynet},
WeakCLIP~\citep{weakclip}) on PlantSeg, plus an ablation that isolates the
impact of binary vs.\ 115-class classification framing. The best pipeline
disease IoU is $33.31\%$, far below the supervised upper bound of
$70.1\%$ (SegNeXt, \citep{segnext}).
\item A reproduction and careful debugging of the recent SPDNet
Siamese plant-disease network~\citep{spdnet} in two fusion regimes
(ADPL-CAM token fusion and a new spatial cross-attention variant), showing
that neither improves pseudo-mask quality over MCTformer under a
classification-only objective.
\item A \emph{shallow-probe localization-capacity methodology} for
measuring how much localization signal a given intermediate layer
\emph{contains}. This is, to our knowledge, the first explicit
upper-bound analysis of this form for plant-disease WSSS features.
Applied to SPDNet, the probe lifts the measurable localization ceiling
from $32\%$ (production cam\_classifier+CRF) to ${\approx}62\%$ disease
IoU, $88\%$ of the fully-supervised SegNeXt baseline.
\item A formal specification and empirical evaluation of three auxiliary
spatial losses (equivariance, patch contrastive, and EMA-teacher
self-distillation) designed to force the Siamese attention map to become
spatially discriminative. We report a clean \emph{null result} with a
root-cause analysis: neither loss injects an independent localization
signal, and we explain why this is a structural property of the chosen
formulations.
\end{enumerate}

\paragraph{Structure of the thesis.}
\Cref{sec:related} reviews the relevant WSSS, foundation-model and
plant-pathology literature. \Cref{sec:setup} describes the data, metrics
and supervised upper bound. \Cref{sec:wsss-baselines} reproduces the
standard WSSS pipeline on PlantSeg. \Cref{sec:spdnet} covers the SPDNet
Siamese approach. \Cref{sec:probe} introduces the shallow-probe
methodology and reports the capacity-analysis results. \Cref{sec:aux}
formalizes the three auxiliary spatial losses and analyses the observed
null result. \Cref{sec:discussion,sec:conclusion} summarize findings and
discuss priorities for remaining work.

% =========================================================
%  2 — Related Work
% =========================================================
\section{Related Work}
\label{sec:related}

\paragraph{Classical WSSS.}
Classical WSSS methods are founded on class activation
maps~\citep{cam,gradcam}, which identify the spatial locations an
image classifier uses to make its decision. Because CAMs tend to be
sparse and focus on the most discriminative parts of an object, a
substantial sub-literature is devoted to \emph{expanding} and
\emph{denoising} these seeds: SEC's seed/expand/constrain framework
\citep{sec}, PSA/AffinityNet pixel-affinity propagation
\citep{affinitynet}, and their Random Walk post-processing. TransCAM
\citep{transcam} and WeakTr~\citep{weaktr} use transformer attention
to refine the raw CAM maps; MCTformer~\citep{mctformer} introduces
class tokens explicitly interacting with patch tokens to produce
class-specific attention maps that are directly usable as CAM surrogates.

\paragraph{Foundation-model-assisted WSSS.}
Large pretrained models such as CLIP~\citep{clip} and SAM~\citep{sam}
provide strong semantic and geometric priors. WeakCLIP~\citep{weakclip}
adapts CLIP to WSSS through a text-to-pixel matching mechanism and
reports PASCAL VOC mIoU in the mid-$60$s. SAM-based WSSS
\citep{sam_pseudo,s2c,fm_wsss} uses SAM as a boundary-precise mask
refiner over CAM-derived prompts. All of these methods make an implicit
\emph{domain assumption}: that the pretrained model's general-purpose
knowledge transfers to the target task. The present work finds that this
assumption is weak for plant diseases, which are niche concepts rarely
represented in web-scale contrastive pretraining.

\paragraph{Plant-disease WSSS and Siamese localization.}
Plant-specific WSSS has been studied in~\citep{plant_wsss} using
CAM\,+\,weak-annotation refinement, and in~\citep{garden,sparsemoesam}
using dedicated plant-lesion architectures with stronger geometric
priors. Most closely related to our Siamese investigation is
SPDNet~\citep{spdnet}, which introduces the ADPL-CAM algorithm. The
SPDNet paper evaluates on only $42$ images with bounding-box IoU and
reports visually compelling but statistically thin results; part of
this thesis is devoted to a careful re-evaluation under pixel-level
IoU on a two-order-of-magnitude larger validation set.

\paragraph{Positioning.}
Our work does not propose a fundamentally new WSSS algorithm. Rather,
we provide a carefully calibrated study of how four families of
approaches — classical CAM pipelines, foundation-model assistance,
Siamese reference-guided CAMs and shallow-probe-based analysis — behave
on a realistic plant-disease benchmark, and introduce a probing tool
that we believe to be a useful unit of analysis for future WSSS work on
specialized domains.

% =========================================================
%  3 — Problem Setting
% =========================================================
\section{Problem Setting, Data and Supervised Upper Bound}
\label{sec:setup}

\subsection{Formal problem statement}

Let $\mathcal{X}\subseteq\mathbb{R}^{H\times W\times 3}$ denote the
image space and let $\mathcal{Y}=\{0,1,\dots,C\}$ denote a label set
consisting of $C$ disease classes and a single background class.
\emph{Fully-supervised} semantic segmentation assumes a training set
$\mathcal{D}_{\mathrm{full}}=\{(x_i,m_i)\}_{i=1}^{N_{\mathrm{f}}}$
where $m_i\in\mathcal{Y}^{H\times W}$ is a pixel-level mask. The WSSS
setting we target assumes only image-level labels:
\[
  \mathcal{D}_{\mathrm{weak}}=\{(x_i,y_i)\}_{i=1}^{N_{\mathrm{w}}},
  \qquad
  y_i\in\{0,1\}^{C},
  \qquad N_{\mathrm{w}}\gg N_{\mathrm{f}},
\]
where $y_{i,c}=1$ indicates the presence of class $c$ anywhere in $x_i$.
The goal is to learn a segmentation function
$f_\theta:\mathcal{X}\to\mathcal{Y}^{H\times W}$ using only
$\mathcal{D}_{\mathrm{weak}}$, evaluated on a held-out set for which
pixel-level masks exist.

\paragraph{Evaluation metrics.}
Given predicted mask $\hat m$ and ground-truth mask $m$, per-class IoU is
$\IoU_c(\hat m,m)=\frac{|\{\hat m=c\}\cap\{m=c\}|}{|\{\hat m=c\}\cup\{m=c\}|}$
and mean IoU is $\mIoU=\frac{1}{C+1}\sum_{c=0}^{C}\IoU_c$. We additionally
report \emph{disease IoU} (the foreground IoU in binary framing), as
$\mIoU$ in a binary task is dominated by the much larger background
class and is misleading as a measure of detection quality.

\subsection{Datasets}

\begin{table}[t]
\centering
\small
\caption{Datasets used in this thesis.}
\label{tab:datasets}
\begin{tabular}{@{}llrll@{}}
\toprule
Dataset & Role & \# Images & Classes & Supervision \\
\midrule
PlantSeg v3 \citep{plantseg}   & Primary, WSSS + eval     & 7\,700   & $1\!+\!115$ & pixel + image-level \\
PlantVillage \citep{plantvillage} & Extra weak labels        & 54\,000  & 38 (mapped) & image-level \\
PASCAL VOC 2012                 & Sanity-check reproduction & 10\,000  & $1\!+\!20$  & pixel + image-level \\
\bottomrule
\end{tabular}
\end{table}

\Cref{tab:datasets} summarizes the three datasets used. PlantSeg contains
$115$ disease classes plus background; PlantVillage contributes
image-level labels mapped to the same $115$-class space via
folder-to-class matching. PASCAL VOC 2012 is used only as a sanity check
that our WSSS pipeline matches published numbers; all reported
plant-domain results are on PlantSeg. The PlantSeg $7\,700/1\,247$
train/val split is respected throughout, and the validation set is never
used at training time, including during refinement stages.

\subsection{Fully-supervised upper bound}

\begin{table}[t]
\centering
\small
\caption{Fully-supervised segmentation on PlantSeg (binary task, 384\,px
input, 30 epochs, cross-entropy loss). SegNeXt
establishes the upper bound.}
\label{tab:upperbound}
\begin{tabular}{@{}lccc@{}}
\toprule
Architecture & $\mIoU$ (\%) & Disease IoU (\%) & BG IoU (\%)\\
\midrule
DeepLabV3+ (ResNet-50)~\citep{deeplabv3p} & 78.9 & 64.6 & 92.9 \\
U-Net (ResNet-34)                         & 78.7 & 64.6 & 92.8 \\
SegFormer (MiT-B3)~\citep{segformer}       & 80.9 & 68.2 & 93.6 \\
\textbf{SegNeXt (MSCAN-T)}~\citep{segnext} & \textbf{82.1} & \textbf{70.1} & \textbf{94.0} \\
\bottomrule
\end{tabular}
\end{table}

\Cref{tab:upperbound} reports the supervised baselines. SegNeXt with the
MSCAN-T backbone reaches $70.1\%$ disease IoU in the binary setting and
serves as the upper bound against which all WSSS numbers are compared. A
multiclass 116-way supervised SegNeXt reaches only $44.2\%$ $\mIoU$
(not tabulated), confirming that the combination of small lesions and
large class cardinality is intrinsically hard even with pixel-level
supervision and motivating the \emph{binary} framing for most of the
WSSS experiments.

% =========================================================
%  4 — Reproduced WSSS Baselines
% =========================================================
\section{Reproduced WSSS Baselines on PlantSeg}
\label{sec:wsss-baselines}

The WSSS pipeline used here is the classical four-stage recipe
illustrated in \Cref{fig:pipeline}. Each stage outputs a pseudo-mask
that serves as the input of the next.

\begin{figure}[t]
\centering
\small
\begin{tabular}{c}
\fbox{\begin{tabular}{c}
\textbf{Image-level labels}\\[2pt]
$\downarrow$ MCTformer classifier\\
\textbf{CAMs}\\
$\downarrow$ DenseCRF (la/ha)\\
\textbf{Refined CAMs}\\
$\downarrow$ PSA (pixel affinity) $+$ Random Walk\\
\textbf{Pseudo-masks}\\
$\downarrow$ WeakCLIP training + CRF\\
\textbf{Final masks}
\end{tabular}}
\end{tabular}
\caption{The standard WSSS pipeline adapted to PlantSeg. Each block is
evaluated independently against the pixel-level ground truth to
diagnose where quality is lost.}
\label{fig:pipeline}
\end{figure}

\subsection{Class activation maps from MCTformer}

\paragraph{MCTformer in one paragraph.}
MCTformer~\citep{mctformer} replaces the single class token of a
standard DeiT~\citep{deit,vit} with $C$ \emph{class tokens},
$\texttt{CLS}_1,\dots,\texttt{CLS}_C$, which interact with patch tokens
through self-attention. For class $c$ and patch $p$, the CAM value is
\begin{equation}
A_c(p) \;=\; \sum_{h=1}^{H_{\mathrm{head}}}
\mathrm{softmax}_p\!\left(\frac{(\texttt{CLS}_c)_h\,(\text{K}_p)_h^{\!\top}}{\sqrt{d_k}}\right),
\label{eq:mctformer_cam}
\end{equation}
averaged over a small set of deep layers. Multi-scale test-time
augmentation (five scales, horizontal flip) is averaged over scales.

\paragraph{Binary vs.\ multi-class framing.}
We train MCTformer in two framings:
\emph{(a)}~\texttt{binary} — a single foreground class ``disease''
vs.\ background; \emph{(b)}~\texttt{MC115} — $115$ disease classes
individually, with the binary CAM obtained post-hoc as
$A_{\mathrm{disease}}(p)=\max_{c=1}^{115} A_c(p)$.

\Cref{tab:wsss_plantseg} summarises the end-to-end disease IoU of the
two framings at each stage of the pipeline. Four observations are
important:

\begin{table}[t]
\centering
\small
\caption{End-to-end WSSS pipeline on PlantSeg (binary evaluation). The
MC115 framing consistently improves disease IoU at every stage. Bold
marks the best per row.}
\label{tab:wsss_plantseg}
\begin{tabular}{@{}lcccc@{}}
\toprule
Stage & Binary $\mIoU$ & Binary $\IoU_{\mathrm{dis}}$ & MC115 $\mIoU$ & MC115 $\IoU_{\mathrm{dis}}$\\
\midrule
LA-CRF         & 45.01 & 12.64 & 51.94 & \textbf{27.77}\\
HA-CRF         & 45.44 & 14.36 & \textbf{53.49} & \textbf{31.16}\\
PSA\,+\,RW     & 43.01 & 20.21 & 42.29 & 30.42\\
WeakCLIP+CRF   & 46.22 & 25.52 & 46.36 & \textbf{33.31}\\
\bottomrule
\multicolumn{5}{l}{\footnotesize Supervised SegNeXt upper bound: $\mIoU=82.1$, $\IoU_{\mathrm{dis}}=70.1$.}
\end{tabular}
\end{table}

\begin{enumerate}[label=(\arabic*),leftmargin=*,itemsep=2pt]
\item The binary MCTformer reaches $99.9\%$ classification mAP because
``is there any disease on this leaf?'' is trivially solvable by
low-level colour/texture features, and consequently produces diffuse,
poorly-localized CAMs.
\item The MC115 framing forces the classifier to distinguish e.g.\
\emph{apple scab} from \emph{apple rust}, which requires attending to
lesion-level features. This structural shift \emph{doubles} the
post-CRF disease IoU ($14.4\%\to31.2\%$).
\item Every subsequent refinement stage (PSA\,+\,RW, WeakCLIP) only
recovers a fraction of the supervised ceiling, and the
WeakCLIP final number is still ${\approx}37$pp below the supervised
upper bound.
\item The global $\mIoU$ and disease $\IoU$ move in opposite
directions across stages: PSA\,+\,RW decreases $\mIoU$ while increasing
disease IoU because it increases recall at the cost of background
precision.
\end{enumerate}

\subsection{Refinement mechanisms tested}

We applied four refinement mechanisms on top of the raw MCTformer CAMs.
\emph{DenseCRF}~\citep{densecrf} with the parameters of
\citep{mctformer} is used twice (low-$\alpha$ and high-$\alpha$ variants).
\emph{PSA\,+\,Random Walk}~\citep{affinitynet} propagates confident
seed regions across a learned pixel affinity graph. \emph{WeakCLIP}
\citep{weakclip} trains a CLIP-based segmentation head on the
pseudo-masks produced above. Finally, we investigated \emph{SAM-1-based
refinement}~\citep{sam,sam_pseudo} with six prompt configurations,
ranging from dense mask prompts to box\,+\,point prompts.

\paragraph{SAM is only useful as a geometric prompt consumer.}
Dense mask prompts (binary or soft) catastrophically collapse to the
leaf boundary ($<\!1\%$ disease IoU). The single configuration that
worked is \emph{box\,+\,CAM-derived points} (experiment~J), which
recovers $33.96\%$ disease IoU — a $+0.65$pp improvement over WeakCLIP
input plus a substantial $+14.92$pp increase in background IoU,
i.e.\ SAM acts as a precision filter rather than as a disease detector.
A point-sampling analysis shows that even at the $99$-th percentile of
CAM activation only ${\sim}53\%$ of the sampled ``positive'' points
actually fall on disease pixels, suggesting that CAM localization
quality itself is the bottleneck rather than the refinement mechanism.

% =========================================================
%  5 — SPDNet
% =========================================================
\section{Siamese Plant-Disease Network (SPDNet)}
\label{sec:spdnet}

\subsection{Architecture}

The Siamese hypothesis of \citep{spdnet} is that pairing a query image
$q$ with a same-class reference $r$ should force the network to focus
on disease-specific features. Our implementation follows the paper's
description while fixing several omissions in the original description
(in particular, an asymmetric application of channel attention and an
unnormalized sum of FPN levels).

Let $\phi$ denote a shared ResNet-50~\citep{resnet} backbone. For input
$x$, $\phi$ produces four feature maps
$\{F^{(k)}\}_{k=1}^{4}$ with channels $256,512,1024,2048$ at progressively
coarser resolutions. An FPN~\citep{fpn} followed by a channel-attention
(MSE) module produces four $256$-channel maps
$P^{(k)}_q$ for the query and $P^{(k)}_r$ for the reference. We consider
two ways of combining them:

\paragraph{Token fusion (ADPL-CAM).}
Each reference level is reduced to a token by global max-pooling,
$T^{(k)}\!=\!\mathrm{GMP}(P^{(k)}_r)\in\mathbb{R}^{256}$, and the fused
query feature is
\begin{equation}
Q_{\mathrm{fused}}(p)
\;=\;\underbrace{\frac{1}{4}\sum_{k=1}^{4}P^{(k)}_q(p)}_{\text{merged query}}
\;+\;\alpha\cdot\frac{1}{4}\sum_{k=1}^{4}T^{(k)},
\label{eq:adpl}
\end{equation}
where $\alpha$ is a learnable scalar initialised at $0.1$.
Crucially, the reference signal becomes a \emph{spatially constant}
bias — the same vector is added at every location.

\paragraph{Spatial cross-attention (SCA).}
To preserve spatial structure, we introduce a multi-head cross-attention
module in which the merged query (of shape $256\times H'\!\times\! W'$)
acts as the query and a $196$-token pooled reference acts as the keys
and values:
\begin{equation}
\mathrm{SCA}(P^{(*)}_q,P^{(*)}_r) \;=\; P^{(*)}_q \;+\;
g\cdot \mathrm{MHA}\!\left(\mathrm{LN}(P^{(*)}_q), \mathrm{LN}(K), \mathrm{LN}(V)\right),
\label{eq:sca}
\end{equation}
with four heads, hidden dimension $256$, and a learnable gate $g$
initialised at $0.1$.

Both fusion modes are fed to a classifier
$\ell(Q_{\mathrm{fused}})=W_{\mathrm{cls}}\cdot\mathrm{GAP}(Q_{\mathrm{fused}})$
with $W_{\mathrm{cls}}\in\mathbb{R}^{115\times256}$. Training uses
multi-label soft-margin loss with AdamW and a cosine schedule with five
warmup epochs at $10^{-4}$ base learning rate.

\subsection{Classification results}

\Cref{tab:spdnet_cls} reports the classification mAP and the post-hoc
disease IoU of the ADPL-CAM-based seeds. Three salient findings stand
out.

\begin{table}[t]
\centering
\small
\caption{SPDNet classification and seed-IoU results. $N$ is the number of
reference images per query. Heavy augmentation = RandAugment + ColorJitter.
Disease IoU is from raw CAMs (no CRF).}
\label{tab:spdnet_cls}
\begin{tabular}{@{}lcccccc@{}}
\toprule
Run & Fusion & $N$ & Aug & val mAP & $\IoU_{\mathrm{dis}}$ (best thr) & Gate $g$\\
\midrule
\texttt{spdnet\_448\_42} (buggy)        & token   & 1 & heavy   & 62.1 & --    & n/a\\
\texttt{spdnet\_fix\_n1\_heavy}          & token   & 1 & heavy   & 85.9 & \textbf{28.1} & n/a\\
\texttt{spdnet\_fix\_n3\_heavy}          & token   & 3 & heavy   & \textbf{89.8} & 24.6 & n/a\\
\texttt{spdnet\_fix\_n3\_light}          & token   & 3 & light   & 89.4 & --   & n/a\\
\texttt{spdnet\_spatial\_n1\_ps}         & spatial & 1 & heavy   & 79.7 & 31.0 & 0.333\\
\texttt{spdnet\_spatial\_n1\_ps\_pv}     & spatial & 1 & heavy   & 88.8 & 28.5 & 0.499\\
\bottomrule
\end{tabular}
\end{table}

\paragraph{Siamese fusion was initially inert.}
In the original implementation, reference features were discarded during
training (fusion was only invoked when \texttt{return\_cam=True}), so
$\alpha$ received zero gradient and the model effectively trained as a
plain ResNet-50\,+\,FPN classifier. Fixing this bug raised val mAP from
$0.621$ to $0.859$ and is a prerequisite for any of the localization
results below.

\paragraph{Higher $N$ helps classification but hurts localization.}
$N{=}3$ references improves val mAP by $+3.9$pp over $N{=}1$ but
\emph{reduces} disease IoU from $28.1\%$ to $24.6\%$, consistent with
the hypothesis that the model over-relies on a global reference bias at
the expense of query spatial structure.

\paragraph{Spatial cross-attention is exercised but does not localize.}
The gate in \Cref{eq:sca} grows from $0.1$ to $0.33/0.50$ during
training, showing that the residual is used. Yet the best disease IoU
from \emph{spatial} fusion (${\sim}31\%$) is no better than token fusion
($28\%$), and $2$pp below MCTformer MC115 (${\sim}31.2\%$ after CRF).
Visualization of attention maps reveals a query-invariant pattern
concentrated on border tokens, rather than a content-driven correspondence
to disease regions — the cross-attention acts as a
\emph{position-conditioned bias} rather than a spatial matcher. This motivates
\Cref{sec:aux}, where we attempt to explicitly shape the attention through
auxiliary losses.

\subsection{Feature-based seeds bypass the classifier projection}

A critical diagnostic experiment was to extract the raw merged query
feature $Q\!=\!\frac14\sum_k P^{(k)}_q$ and turn it into a 1-channel
seed map by simple channel aggregation
\emph{before} the classifier projection:
\begin{align}
\mathrm{feat\_chmean}(Q)(p) &\;=\; \tfrac{1}{256}\textstyle\sum_{c=1}^{256} Q_c(p),\\
\mathrm{feat\_chvar}(Q)(p)  &\;=\; \tfrac{1}{256}\textstyle\sum_{c=1}^{256}\!\left(Q_c(p)-\bar Q(p)\right)^2.
\end{align}
The resulting maps, thresholded and CRF-refined with a per-distribution
parameter sweep, yield disease IoU of $36.5\%\to42.1\%$ (feat\_chmean
$\to$ feat\_chmean+CRF with $\sigma_{\mathrm{rgb}}{=}5$,
$\tau_{\mathrm{bg}}{=}0.30$). This $+6.5$pp improvement over MCTformer
and $+14.1$pp over ADPL-CAM is achieved \emph{without retraining} and
strongly suggests that the classifier-head projection $W_{\mathrm{cls}}$
destroys useful localization information present in the backbone
features — an observation that directly motivates the shallow-probe
methodology of \Cref{sec:probe}.

% =========================================================
%  6 — Shallow Probe
% =========================================================
\section{Localization-Capacity Analysis via Shallow Probes}
\label{sec:probe}

The previous section showed that \emph{more} localization information
is present in the backbone than in the classifier-projected CAMs. We now
quantify this systematically: for any intermediate representation $\phi$
of the network, we want to know the \emph{upper bound} on the
localization information that $\phi$ can provide, independent of any
particular aggregation heuristic.

\subsection{The shallow-probe argument}

Let $\phi_\theta:\mathcal{X}\to\mathbb{R}^{d\times H'\!\times\! W'}$ be
a layer of a trained network with frozen parameters $\theta$. Let
$\mathcal{A}$ denote the class of all \emph{unsupervised} aggregations
from $\mathbb{R}^{d}$ to $[0,1]$ — i.e.\ any hand-designed function
$a:\mathbb{R}^d\!\to\![0,1]$ applied pointwise, such as channel-mean,
max, GradCAM, or feat\_chvar. The localization IoU any such aggregator
can provide on ground-truth mask $m$ is at most
\begin{equation}
\sup_{a\in\mathcal{A}}\;\IoU\!\bigl(a(\phi_\theta(x)),\,m\bigr),
\label{eq:unsup_bound}
\end{equation}
which is hard to evaluate because $\mathcal{A}$ is infinite. Instead,
consider a \emph{shallow supervised probe} $g_\psi$ with a limited
parameter budget $B$ (e.g.\ two $1{\times}1$ convolutions with 64 hidden
channels, $\approx\!17{,}000$ parameters). Training $g_\psi$ with frozen
$\phi_\theta$ on pixel-level supervision gives
\begin{equation}
\IoU^{*}_{\mathrm{probe}}(\phi_\theta) \;\coloneqq\;
\max_{\psi}\;\mathbb{E}_{(x,m)}\!\bigl[\IoU\!\bigl(g_\psi(\phi_\theta(x)),\,m\bigr)\bigr].
\label{eq:probe_def}
\end{equation}
Because the probe has direct access to pixel labels and $g_\psi$ can
realize any linear (and mild non-linear) combination of feature
channels, the probe is \emph{at least} as good as any unsupervised
$a\in\mathcal{A}$ on average — up to the probe's own capacity limits
and to estimation noise. When the probe capacity is small enough for
$g_\psi$ to be the bottleneck, the probe's IoU is dominated by the
informativeness of $\phi_\theta$, which we denote
$I(\phi_\theta\!\to\! m)$:
\begin{equation}
\IoU^{*}_{\mathrm{probe}}(\phi_\theta)
\;\approx\;\tilde I(\phi_\theta\!\to\! m)
\;\geq\;\sup_{a\in\mathcal{A}}\IoU\!\bigl(a(\phi_\theta(x)),\,m\bigr).
\end{equation}
Thus the probe is a \emph{tractable empirical upper bound} on the
unsupervised localization capacity of a specific feature map. The
critical caveats are:
(i)~$g_\psi$ must be shallow enough that its capacity is not
itself the bottleneck, otherwise the bound is not tight;
(ii)~the bound is evaluated on a held-out set and relies on the
standard generalization-gap argument of empirical risk minimization.

\subsection{Three-phase probe protocol}

\begin{table}[t]
\centering
\small
\caption{Shallow probe capacity analysis on SPDNet. ``S'' is the
composite score $\max(\text{probe},\text{chmean},\text{chvar},\text{cam\_cls})$
after CRF. All numbers are CRF-refined disease IoU (\%) on PlantSeg val.
Ph1 and Ph2 use a deterministic 300-image subset; Ph3 uses the full
1247 val set.}
\label{tab:probe}
\begin{tabular}{@{}llcccccc@{}}
\toprule
Phase & Ckpt & Position & Probe & chmean & chvar & cam\_cls & $S$\\
\midrule
Ph1 (frozen) & token     & P1 layer4              & 46.2 & 35.1 & 34.6 & --   & 46.2\\
             & token     & P3 query\_merged       & 43.2 & 36.9 & 37.0 & 32.4 & 43.2\\
             & spatial   & P3 query\_merged       & 42.6 & 22.0 & 36.9 & 27.6 & 42.6\\
\midrule
Ph2 (unfrozen)&token     & P3 query\_merged       & 59.8 & 33.8 & 39.8 & 38.3 & 59.8\\
             & spatial   & \textbf{P3 query\_merged}&\textbf{59.9}& 22.0 & 40.4 & 37.0 & \textbf{59.9}\\
             & spatial   & P4 fused               & 59.0 & 21.9 & 37.1 & 29.9 & 59.0\\
\midrule
Ph3 (from scratch) & spatial & P3 query\_merged   & \textbf{61.8} & 21.9 & 24.3 & 30.8 & \textbf{61.8}\\
\bottomrule
\end{tabular}
\end{table}

We instantiate the argument in three phases of increasing training
freedom, all with the same probe architecture
($\mathrm{Conv}_{1\!\times\!1}(d,64)\!\to\!\mathrm{BN}\!\to\!\mathrm{ReLU}\!\to\!\mathrm{Conv}_{1\!\times\!1}(64,1)$)
and a composite BCE\,+\,Dice segmentation loss:

\begin{enumerate}[leftmargin=*,itemsep=1pt]
\item \textbf{Phase 1 — Frozen.} $\theta$ is frozen; only $\psi$ is
updated for 20 epochs. This is the cleanest form of the upper bound in
\Cref{eq:probe_def}: it measures what the current feature map contains
without letting the backbone adapt.
\item \textbf{Phase 2 — Unfrozen joint.} $\theta$ and $\psi$ are both
trained for 15 epochs with a segmentation + classification composite
loss $\lambda_{\mathrm{seg}}\mathcal{L}_{\mathrm{seg}}
+\lambda_{\mathrm{cls}}\Lcls$, both weights $=1$. This measures the
capacity of the \emph{architecture}, not just the current features.
\item \textbf{Phase 3 — From scratch.} The entire backbone plus the
probe is trained from a random initialization under the same joint loss.
This is the from-scratch ceiling of the chosen architecture under pixel
supervision.
\end{enumerate}

\subsection{Results}

\Cref{tab:probe} reports the three-phase disease-IoU numbers for the
SPDNet spatial checkpoint trained on PlantSeg\,+\,PlantVillage, plus a
comparable token-fusion checkpoint. The key findings are:

\begin{enumerate}[label=(\alph*),leftmargin=*,itemsep=2pt]
\item \textbf{Frozen capacity is high.} Even with a frozen backbone, a
two-layer probe reaches $46.2\%$ disease IoU at P1\_layer4 and $42.6\%$
at P3\_query\_merged — $14$pp above the best unsupervised aggregation
(\texttt{chvar}, $36.9\%$). The localization signal is therefore
\emph{present} in the feature space; it is not being extracted by any
hand-crafted aggregation.
\item \textbf{Fine-tuning narrows probe-position differences.} After
15 epochs of joint training (Phase 2), the spread between probe
positions shrinks from $24$pp (Phase 1) to $4$pp (Phase 2): the
choice of probe position becomes nearly irrelevant because the
backbone has adapted to the seg objective at all positions.
\item \textbf{Upper bound $\approx62\%$.} Phase 3 reaches
$61.8\%$ after CRF and $64.9\%$ with raw thresholding — only
$8$pp below the fully-supervised SegNeXt benchmark of $70.1\%$ and
\emph{$29$pp above} the production WSSS pipeline. This defines a
new, much-higher \emph{empirical upper bound} for WSSS on this
architecture class.
\item \textbf{Spatial vs.\ token is a closed question.} In every
measurement (classification mAP, cross-attention viz, probe IoU),
spatial fusion and token fusion are indistinguishable after
fine-tuning. Section~\ref{sec:spdnet}'s conclusion — that a spatial
cross-attention alone is insufficient — is robustly confirmed.
\item \textbf{P6\_attn\_map loses to chmean.} When the
1-channel spatial attention map is used as a seed, the probe reaches
only $17$--$22\%$ disease IoU, i.e.\ the attention map as produced by a
classification-only objective is \emph{less} informative than a naive
channel-mean baseline.
\end{enumerate}

The practical implication is that the current WSSS pipeline leaves
${\approx}30$pp of disease IoU on the table. Replacing the
\texttt{cam\_classifier}+CRF seed source with a probe-generated seed
source (Phase 2 recipe) is expected to lift the end-to-end pipeline
accordingly, and constitutes the single highest-leverage planned
follow-up experiment.

% =========================================================
%  7 — Auxiliary Spatial Losses
% =========================================================
\section{Auxiliary Spatial Losses: Formalization and Null Result}
\label{sec:aux}

The shallow-probe analysis identifies \emph{what is possible} but the
best \emph{fully-unsupervised} seed is still the feature-based
\texttt{chmean}/\texttt{chvar} aggregation at $\approx\!37\%$. To close
this gap while remaining within the WSSS paradigm, we tested whether
auxiliary losses applied during SPDNet training can force the network
to learn a spatially discriminative attention map, without any
pixel-level labels.

\subsection{Notation}

For a batch of $B$ query/reference pairs $(q_i,r_i)$:
$F^{P3}_q\!\in\!\mathbb{R}^{B\times 256\times H'\!\times\!W'}$ denotes
the pre-fusion merged-query feature (``Point P3''), and
$F^{P4}\!=\!\mathrm{SCA}(F^{P3}_q,F^{P3}_r)$ the post-fusion feature
(``Point P4''). Let $M(q,r)\!\in\!\mathbb{R}^{B\times H'\!\times W'}$
be a \emph{single-channel attention-concentration map}, and
$S\!=\!W_{\mathrm{cls}}F^{P4}\!\in\!\mathbb{R}^{B\times C\times H'\!\times W'}$
the per-class spatial logits. Let $\mathcal{T}\!=\!\{\mathrm{id},
\mathrm{hflip},\mathrm{rot90},\mathrm{rot180},\mathrm{rot270}\}$ be
five discrete, invertible geometric transforms.

\paragraph{Attention-concentration map.}
A first implementation used $M(q)\!=\!\mathbb{E}_{k}[\mathrm{attn}_{q,k}]$,
but because post-softmax rows sum to $1$ this equals the constant
$1/N_{\mathrm{ref}}$ identically, making any equivariance loss vacuous.
The corrected map is the row-wise normalized negentropy,
\begin{equation}
M(q_i) \;=\; 1 + \frac{1}{\log N_{\mathrm{ref}}}
\sum_{k=1}^{N_{\mathrm{ref}}} p_{i,k}\log p_{i,k}\;\in\;[0,1],
\label{eq:attn_conc}
\end{equation}
where $p_{i,k}$ is the softmax attention weight from query location $q_i$
to reference key $k$. $M=0$ corresponds to uniform attention,
$M=1$ to perfect peaking.

\subsection{Equivariance loss $\Leq$}

A model that truly localizes must produce attention that tracks the
query's spatial content. Given a transform $T\!\in\!\mathcal{T}$ sampled
uniformly per batch and applied only on the query branch:
\begin{equation}
\Leq(q,r) \;=\; \frac{1}{B H'W'}\sum_{i=1}^{B}
\bigl\|\,T\!\left(M(q_i,r_i)\right)-M\!\left(T(q_i),r_i\right)\bigr\|_{F}^{2}.
\label{eq:Leq}
\end{equation}

\subsection{Patch contrastive loss $\Lcon$}

A projection head $g:\mathbb{R}^{256}\!\to\!\mathbb{R}^{128}$ (single
$1{\times}1$ convolution, $L^2$-normalized) embeds the pre-fusion
query feature $F^{P3}_q$ into $z_{i,p}\!\in\!S^{127}$. For each image
$i$ with positive class $y_i$ we compute a bootstrap per-class spatial
map
$\tilde S_{i,y_i}\!=\!\mathrm{minmax}(W_{\mathrm{cls},y_i}F^{P4}_i)$
and define, under \texttt{torch.no\_grad}:
\begin{align}
\mathcal{A}_i &= \mathrm{topK}_p\,\tilde S_{i,y_i}(p),\quad K=8, &
\mathcal{B}_i &= \{p:\tilde S_{i,y_i}(p)<\text{median}\},\,|\mathcal{B}_i|=M=16.
\end{align}
Positive peers $\mathcal{P}(a)=\mathcal{A}_i\!\setminus\!\{a\}$ and
negative peers $\mathcal{N}(a)=\mathcal{B}_i\cup\bigcup_{j:y_j\neq y_i}\mathcal{A}_j$.
The SupCon-style contrastive loss~\citep{simclr} is
\begin{equation}
\Lcon = \frac{1}{\sum_i|\mathcal{A}_i|}
\sum_{i}\sum_{a\in\mathcal{A}_i}
\!\!-\frac{1}{|\mathcal{P}(a)|}\!\!\sum_{p\in\mathcal{P}(a)}
\!\!\log\frac{\exp(z_a^{\!\top} z_p/\tau)}
{\sum_{x\in\{p\}\cup\mathcal{N}(a)}\exp(z_a^{\!\top} z_x/\tau)},
\label{eq:Lcon}
\end{equation}
with $\tau=0.07$. Anchors are taken from $F^{P3}$ (pre-fusion) so that
the loss cannot ``cheat'' by aligning the post-fusion residual.

\subsection{Self-distillation $\Ldist$ (DINO-style)}

A teacher $\theta_t$ is an EMA copy of the student $\theta_s$ with
momentum $\alpha=0.999$ (buffers included). For image $i$ with
positive class $y_i$, $\tilde S^{(\cdot)}_i\!=\!\mathrm{flatten}(S^{(\cdot)}_{i,y_i})$.
With centering update $c\leftarrow\beta c+(1-\beta)\frac1B\sum_i\tilde S^{(t)}_i$
and temperatures $T_t{=}0.04,T_s{=}0.1$:
\begin{equation}
P^{(t)}_i\!=\!\mathrm{softmax}\!\left(\!\tfrac{\tilde S^{(t)}_i-c}{T_t}\!\right)\!,\quad
P^{(s)}_i\!=\!\mathrm{softmax}\!\left(\!\tfrac{\tilde S^{(s)}_i}{T_s}\!\right)\!,\quad
\Ldist\!=\!\frac{1}{B}\!\sum_i\!\mathrm{KL}(P^{(t)}_i\,\|\,P^{(s)}_i).
\label{eq:Ldist}
\end{equation}
Centering mitigates mode collapse, temperature sharpening mitigates the
trivial-copy fixed point, and a warm-up of $10$ epochs
($\lambda_{\mathrm{distill}}{=}0$) avoids cold-start pathologies.

\subsection{Composite objective and experiments}

\begin{equation}
\Ltotal \;=\; \Lcls \;+\; \lambda_{\mathrm{eq}}\Leq \;+\;
\lambda_{\mathrm{con}}\Lcon \;+\; \lambda_{\mathrm{distill}}\Ldist.
\end{equation}
Launch defaults: $\lambda_{\mathrm{eq}}\!=\!1.0$, $\lambda_{\mathrm{con}}\!=\!0.5$,
$\lambda_{\mathrm{distill}}\!=\!0$. Three runs were carried out:
\emph{baseline} (eq-only, $80$\,ep, from scratch); \emph{C} (warmstart
from the eq-only checkpoint with $\Lcon$ turned on for $40$\,ep);
\emph{F} (from scratch with $\Lcon$ linearly ramped from
epoch~$14$ to $21$, reaching $\lambda=0.5$).

\begin{table}[t]
\centering
\small
\caption{Auxiliary-loss experiments on SPDNet (spatial fusion,
PlantSeg+PlantVillage). Online CAM-IoU is measured on a deterministic
100-image val subset every two epochs. ``History std'' is the
epoch-level standard deviation.}
\label{tab:aux}
\begin{tabular}{@{}lcccccc@{}}
\toprule
Run & Epochs & Final mAP & Best CAM-IoU & $\Delta$ CAM-IoU & History std & $\Lcon$ change\\
\midrule
Baseline ($\Leq$) & 80 & 0.843 & 0.246 & $+0.048$ & 0.016 & --\\
Warmstart C       & 40 & 0.841 & 0.252 & $+0.002$ & 0.006 & $0.028\!\to\!0.003$\\
Warmup F          & 62 & 0.768 & 0.234 & $+0.036$ & 0.011 & similar drop\\
\bottomrule
\end{tabular}
\end{table}

\subsection{Null result and root-cause analysis}

\Cref{tab:aux} tells a clean story. Activating $\Lcon$ at strength
$\lambda{=}0.5$ (run~C) reduces the loss value by a factor of $10$
while moving the quantity it is supposed to improve
(\texttt{val/cam\_iou\_best}) by only $+0.002$, less than one standard
deviation of its own history. Run~F shows the same pattern in reverse:
during the ramp, the CAM-IoU metric \emph{drops} from $0.246$ to
$0.219$. Independently, Phase-1 probing on the baseline eq-only
checkpoint reveals probe-IoUs indistinguishable within $\pm 1$pp from
the classifier-only checkpoint of \Cref{sec:probe}.

Two structural failure modes explain the result:

\paragraph{(F1) $\Lcon$ anchors are classifier-bootstrapped.}
Inspection of $\Lcon$ reveals that the top-$K$ anchor set $\mathcal{A}_i$
in \Cref{eq:Lcon} is defined by the current student classifier's own
per-class argmax. The loss therefore pulls together embeddings where
\emph{the current classifier already thinks} the disease lives, and
pushes apart embeddings where it doesn't. Any classifier with
internally-consistent (but possibly wrong) spatial beliefs minimizes
$\Lcon$ — this is a self-distillation of the classifier's current
beliefs, dressed up as a contrastive loss, and carries no independent
localization signal. This is precisely what the $10\times$ drop of
$\Lcon$ with unchanged CAM-IoU demonstrates empirically.

\paragraph{(F2) $\Leq$ has nothing to preserve.}
The Phase-1 probe of P6\_attn\_map shows that the attention map as
learned by classification-only training is less informative than
channel-mean aggregation, i.e.\ it is functionally near-uniform. For
any transform $T$, $T(\text{const})=\text{const}$, so $\Leq$ is
trivially minimized by the existing uniform map. $\Leq$ can only
\emph{preserve} existing spatial structure; it cannot create it. Since
classification does not need attention structure, classification does
not create any, and $\Leq$ has nothing to act on.

\paragraph{Redesign priorities.}
The analysis above suggests that \emph{any} self-referential anchor
source ($\Lcon$) will fail. A minimum follow-up experiment that breaks
circularity is to replace the student-argmax anchor by an EMA-teacher
anchor (the \texttt{EMATeacher} class used for $\Ldist$ is already
available and tested). Orthogonally, a \emph{concentration regularizer}
$\mathbb{E}[H(\mathrm{attn})]$ or a $\mathrm{TV}$-smoothness-plus-mass
term would give $\Leq$ structure to preserve. The strongest proposed
follow-up, however, is to distill the Phase-2 probe head's prediction
into the SCA attention map via a masked MSE or KL, providing the
explicit localization signal that $\Leq/\Lcon$ fail to inject; the
probe teacher already reaches $60\%$ disease IoU on the same
architecture family.

% =========================================================
%  8 — Discussion
% =========================================================
\section{Discussion}
\label{sec:discussion}

\paragraph{Why is plant-disease WSSS so much harder than PASCAL VOC?}
Our pipeline achieves $64.0\%$ $\mIoU$ on PASCAL VOC (WeakCLIP on
VOC pseudo-masks) and only $46\%$ $\mIoU$ with $33\%$ disease IoU on
PlantSeg, using the same code and hyperparameters. Three structural
factors contribute. First, PASCAL VOC's classes are large,
well-localized objects whose most discriminative regions coincide with
the object itself; plant-disease lesions are often small, near the
boundaries of the leaf and low-contrast against healthy tissue, so the
discriminative region for classification
(``the image contains an apple leaf'') does not localize the disease.
Second, PlantVillage is not a drop-in replacement for the PASCAL label
noise distribution — it is photographed on uniform backgrounds that
make disease recognition trivial, providing near-perfect
classification mAP but little useful spatial signal. Third, the
pretraining objectives of ViT, CLIP and SAM are dominated by natural
images with much broader semantic distribution than the
niche plant pathology categories.

\paragraph{Why did SPDNet's Siamese structure fail to help?}
Three independent diagnostics (classification mAP in \Cref{tab:spdnet_cls},
cross-attention visualizations, and probe signatures in
\Cref{tab:probe}) converge on the same conclusion: a classification-only
training objective does not provide a learning signal that forces the
attention to be spatially discriminative. This is consistent with the
classical observation in WSSS~\citep{sec,affinitynet} that
classification networks focus on most-discriminative parts rather than
on the full object. The SPDNet paper's $42$-image evaluation
\citep{spdnet} with bounding-box IoU is not representative of the
behavior seen at PlantSeg's scale.

\paragraph{How large is the information-theoretic gap that remains?}
The shallow-probe upper bound of $62\%$ disease IoU versus the
supervised baseline of $70\%$ suggests that ${\approx}8$pp of the gap is
\emph{not} a property of the backbone features; it is either the CRF
over-smoothing of already-strong seeds (empirically measured at
${\approx}3$pp in \Cref{tab:probe}), the capacity limit of the
two-layer probe head itself, or an architectural property of ResNet-50.
The latter would require a bigger head, multi-scale outputs, or a
different backbone, and is out of scope for the present thesis.

\paragraph{Threats to validity.}
The results rely on the standard train/val split of PlantSeg,
without a test set. Reported CRF parameters were each tuned on
distribution-matched sweep subsets to avoid the optimistic bias of
reusing VOC-tuned parameters, but the CRF sweep itself consumes val
supervision and so the reported CRF-refined numbers are mildly
optimistic. The probe methodology gives an upper bound on
\emph{unsupervised} localization capacity; WSSS pseudo-mask training
cannot exploit this directly without an intermediate stage. All
reported results are preliminary in the sense that additional
ablations, a probe-seeded end-to-end WSSS pipeline and a
non-circular-anchor re-run of $\Lcon$ are in progress.

% =========================================================
%  9 — Conclusion
% =========================================================
\section{Conclusion and Future Work}
\label{sec:conclusion}

This thesis reported an end-to-end experimental study of weakly-supervised
semantic segmentation for plant-disease localization on the PlantSeg
dataset. Starting from a supervised upper bound of $70.1\%$ disease
IoU (SegNeXt), we reproduced three established WSSS baselines
(MCTformer, PSA/Random Walk, WeakCLIP), reached a best end-to-end
disease IoU of $33.3\%$, and showed that the Siamese SPDNet architecture
in two fusion regimes does not close this gap under a
classification-only objective. Two methodological contributions
followed: (i)~a shallow-probe protocol that measures an empirical
upper bound on the localization information contained in any
intermediate layer, which revealed that $\approx\!62\%$ disease IoU is
already latent in SPDNet features but is destroyed by the final
classifier projection; (ii)~a formal specification and clean null
result for two commonly-proposed auxiliary spatial losses
($\Leq$, $\Lcon$), traced to structural properties of their
formulations (self-referential anchors and a near-uniform starting
attention map, respectively).

\paragraph{Priorities for remaining work.}
\begin{enumerate}[leftmargin=*,itemsep=1pt]
\item \textbf{Promote probe seeds to production.} Replace the current
\texttt{cam\_classifier}+CRF seed source of Steps~1--2 of the WSSS
pipeline with the Phase-2 probe-head recipe
(\Cref{sec:probe}); the expected gain is $20$--$25$pp in disease IoU on the
final WeakCLIP masks.
\item \textbf{Re-tune CRF on strong seeds.} The current CRF
parameters, tuned on weaker classifier-only seeds, over-smooth
Phase-3-quality inputs (\Cref{tab:probe}). A fresh parameter sweep on
$\geq\!200$ images is low-effort and expected to recover $\approx\!3$pp.
\item \textbf{Non-circular $\Lcon$.} Re-run the C/F experiments with
EMA-teacher anchors (\texttt{EMATeacher} class) instead of student-argmax
anchors. Success is defined as \texttt{val/cam\_iou\_best} moving above
the eq-only plateau by at least two history standard deviations
($\geq\!+0.03$).
\item \textbf{Probe-to-attention distillation.} Train the SCA
attention map with a masked-MSE or KL distillation from the Phase-2
probe head's prediction — an explicit localization signal that sits
outside the WSSS-premise-violating regime of using ground-truth masks at
training time.
\end{enumerate}

All source code, DVC-tracked artifacts, MLflow runs and reports are
publicly available at
\url{https://github.com/dfbakin/plant-diseases-segmentation}, enabling
independent reproduction and extension of every experimental number
reported in this document.

% =========================================================
%  Bibliography
% =========================================================
\bibliographystyle{unsrtnat}
\bibliography{project_proposal_Bakin}

\end{document}
